\section{Conclusion}
\label{sec:conclusion}

Reading a clinical concept from a vision-language model's visual representation does not make its probe direction a handle on that concept. Across \cfCheckpoints{} checkpoints and two chest-radiograph datasets, concepts are readable in \cfChestReadablePct\% and answerable in \cfChestAnswerablePct\% of cells but owned in \cfChestOwnedPct\%. The same procedure owns \cfCocoOwnedPct\% of natural-object cells. Non-clinical attributes of the same films are owned in \cfAttrAllAttrOwned{} of \cfAttrAllAttrN{} cells against \cfAttrAllClinOwned{} of \cfAttrAllClinN{} for findings, and in none of the \cfStratAttrTopN{} attribute cells the model answers well: the deficit follows the chest-radiograph representation and its reader, not the concept's clinical status. The model's answer direction is a handle. On chest radiographs, it is nearly orthogonal to the probe direction in raw model space (cosine $\cfAnsCosChestMed$) and far from it under the training-feature covariance (whitened cosine $\cfValCosSigmaChestMin$--$\cfValCosSigmaChestMax$), compared with $\cfAnsCosCocoMed$ and $\cfValCosSigmaCoco$ on natural images. Checkpoints sharing a vision tower differ in what they own. Crossing the tower with the reader leaves every clinical concept unwritable in all four combinations. Ownership therefore rests on how the reader uses the representation.